% RESULTADO_RECUPERADO: c40_termicas.tex from the active integral.
% Extended proof: mode counting, thermal moments, spectral maxima,
% and substitution of the calendar scale.
% The realization hypotheses remain explicit.
\section{Thermal radiation and the information scale of the calendar}
\label{v:sec:radiacion-termica}

Occupation statistics provide the link between the modes of
a realization and their mean energy. In the electromagnetic realization,
the dispersion relation and polarization multiplicity additionally determine the
spectral density. Their composition yields the thermal moments
and spectral maxima, while keeping three operations distinct:
production of the action scale, representation of the modes, and
evaluation of the thermal state.

\subsection{Action, calendar, and temperature}

Retain the clock orientation \(a\) and abbreviate
\(\hbar=\hbar_a\), \(c=c_{\rm int}=c_{\rm vac}\), with the
identification in \eqref{v:eq:identidad-velocidad} made in the same chart.
The energy \(\varepsilon_{108}\) in this section is
\(\epsilon_a\) of \eqref{v:ant:eq:ii-kb-aut-energia};
\(\Theta_{108}=\Theta_{\rm clk,a}\). Dividing it by \(\log3\)
gives the energy per nat in \eqref{v:ant:eq:ii-kb-aut-energia-nat}.
The construction and covariance of the thermal pair are given in
Section~\ref{v:ant:sec:ii-kb-aut-desarrollo}.

Write \(h=2\pi\hbar>0\) for the action representation, \(t_0>0\)
for the temporal unit of its dimensional chart, and \(\ell_0=ct_0\) for
the corresponding length. The chronological cell of length \(108\)
determines the energy
\begin{equation}
 \varepsilon_{108}=\frac{h}{108t_0}
 =\frac{\hbar c}{L_\alpha}=Q_{L,a}=E_{P,a},\qquad
 k_B\Theta_{108}\log 3=Q_{L,a},\qquad
 \beta_{108}=\frac{108t_0\log3}{h}.
 \label{v:eq:energia-informacion-calendario}
\end{equation}
Here \(\beta=(k_B\Theta)^{-1}\) is the Gibbs parameter. The equality
fixes the energy-valued product of temperature and the entropy
transducer. The dimensional specification of \(k_B\) and
\(\Theta_{108}\) separately uses the thermal section and its bases. The natural
logarithm measures information per tritic branch; its factor \(\log3\) converts
the cardinality of the cell into dimensionless entropy.

Also introduce the associated units
\begin{equation}
 t_P=\frac{108}{2\pi_{\HMT}}t_0=\frac{L_\alpha}{c},\qquad
 \ell_P=ct_P=L_\alpha,\qquad
 E_P:=E_{P,a}=\frac{\hbar}{t_P}=Q_{L,a}=\varepsilon_{108}.
 \label{v:eq:unidades-termicas-calendario}
\end{equation}
These equalities specify the meaning of the symbols in this section;
their evaluation in another chart preserves the identities between quotients.
The length \(\ell_P=(54/\pi_{\HMT})\ell_0\) first coincides with the
structural length \(L_\alpha\) of \eqref{v:eq:longitud-alpha}; the
subsequent gravitational identity \(\ell_P^2=G_a\hbar_a/c^3\) is the
corollary \eqref{v:eq:gravedad-unica-alpha}, not the origin of that scale.

The centred equivalence
\eqref{v:eq:equivalencia-termica-centrada} and the transducer identity
\eqref{v:eq:identidad-transductor-numeral} fix the energy scale and its
thermal reading before the electromagnetic realisation. The marker
$p_0$ in \eqref{v:eq:particion-portador-termico} retains normalisation
and origin when a branch is conditioned. This substitution determines
$\beta$ and the energy scale; it does not retrospectively select the
dispersion relation or the two polarisations specified below. The
Planck, Stefan--Boltzmann, and Wien laws therefore remain downstream
consequences of the declared modal realisation.

\subsection{Mode density and spectral law}

\begin{proposition}[Thermal realization of transverse modes]
\label{v:prop:ley-espectral}
Consider a free realization in a three-dimensional periodic box,
whose excitation space contains the wave vectors
\(\mathbf k\ne0\), with two transverse modes per vector, dispersion
\(\omega=c|\mathbf k|\), excitation energy \(\hbar\omega\),
zero chemical potential, and \(\beta>0\). The constant mode is kept
fixed as background data. In the thermodynamic limit, the energy density
per angular frequency is
\begin{equation}
 u_\omega(\Theta)=\frac{\hbar\omega^3}{\pi^2c^3}
 \frac{1}{e^{\beta\hbar\omega}-1},\qquad \omega>0.
 \label{v:eq:planck-angular}
\end{equation}
The energy is that of the excitations above the vacuum of this
representation.
\end{proposition}

\begin{proof}
For a box of volume \(V\), the wave vectors have density
\(V/(2\pi)^3\). In a radial shell of thickness \(dk\), the density
per unit volume, including the two polarizations, is
\[
 2\frac{4\pi k^2\,dk}{(2\pi)^3}=\frac{k^2}{\pi^2}\,dk
 =\frac{\omega^2}{\pi^2c^3}\,d\omega.
\]
This is the limiting density of the sums over the periodic lattice
with the constant mode omitted. This choice defines the partition product
before taking the limit: each excitation mode has positive energy.
Excluding a single point preserves the limiting radial density.
The previously obtained bosonic occupation is
\(\bar n_\omega=(e^{\beta\hbar\omega}-1)^{-1}\).
Multiplying the density by \(\hbar\omega\bar n_\omega\) proves
\eqref{v:eq:planck-angular}. Near zero, the energy integrand is
of order \(\omega^2\), and at infinity it decays exponentially, so
its integral converges.
\end{proof}

The density arises from the specified three-dimensional representation.
The transverse multiplicity used here has that physical domain;
the dimension of a harmonic space on the APP torus describes,
separately, the cohomology of its cellular complex.

% FORMALIZACION_NUEVA of the justification for the limit used in
% 70_radiacion_termica.tex. It does not introduce a constant generator.
\subsection{Thermodynamic limit of occupation sums}
\label{v:subsec:limite-termodinamico}

The preceding radial density can be obtained as an effective limit of the
sums in the periodic realization. The estimate must retain the fixed
constant mode and uniformly control both the origin and the spectral
tail. The action and calendar representations remain prior
to this realization operation.

\begin{lemma}[Convergence of the three thermal observables]
\label{v:lem:limite-termodinamico}
Let \(a,b>0\), \(L>0\), \(\Delta=2\pi/L\), and
\[
 F_N(r)=\frac{1}{e^{ar}-1},\qquad
 F_E(r)=\frac{br}{e^{ar}-1},\qquad
 F_Z(r)=-\log(1-e^{-ar})\quad(r>0).
\]
For each of these functions,
\begin{equation}
 \lim_{L\to\infty}\frac{2}{L^3}
 \sum_{\nu\in\mathbb Z^3\setminus\{0\}}
 F\!\left(\frac{2\pi}{L}|\nu|\right)
 =\frac{1}{\pi^2}\int_0^\infty r^2F(r)\,dr.
 \label{v:eq:limite-modos-discretos}
\end{equation}
The sums and the integral are finite. The specialization
\(a=\beta\hbar c\), \(b=\hbar c\) gives, respectively,
the number density, excitation-energy density, and logarithmic partition-function
density of the transverse modes.
\end{lemma}

\begin{proof}
All three functions are continuous and positive on \((0,\infty)\).
There exist constants \(C_0,C_1,d>0\), depending on \(a,b\) but
independent of \(\Delta\), such that
\begin{equation}
 F(r)\leq\frac{C_0}{r}\quad(0<r\leq1),\qquad
 F(r)\leq C_1e^{-dr}\quad(r\geq1).
 \label{v:eq:envolventes-termicas}
\end{equation}
Indeed, \(e^{ar}-1\geq ar\) bounds \(F_N\) and \(F_E\)
near the origin; moreover,
\(F_Z(r)=\log(1+F_N(r))\leq F_N(r)\leq1/(ar)\).
For \(r\geq1\), the inequality
\(-\log(1-y)\leq y/(1-y)\), with \(y=e^{-ar}\), gives the
exponential bound for \(F_Z\). The other two follow from
\((e^{ar}-1)^{-1}\leq e^{-ar}/(1-e^{-a})\), absorbing the
factor \(r\) in \(F_E\) into an exponential of smaller decay rate. These bounds
already prove integrability of \(r^2F(r)\).

For the discrete sums, group the lattice points by
\(m=\|\nu\|_\infty\geq1\). Each shell contains exactly
\[
 (2m+1)^3-(2m-1)^3=24m^2+2
\]
points and satisfies \(m\leq|\nu|\leq\sqrt3m\).
Fix \(0<\delta\leq1\) and \(0<\Delta\leq1\).
If \(\delta<\Delta\), there are no points with
\(0<\Delta|\nu|\leq\delta\). Otherwise, setting
\(N=\lfloor\delta/\Delta\rfloor\), the first bound in
\eqref{v:eq:envolventes-termicas} gives
\begin{align*}
 \Delta^3\!\sum_{0<\Delta|\nu|\leq\delta}F(\Delta|\nu|)
 &\leq C_0\Delta^2\sum_{m=1}^{N}\frac{24m^2+2}{m}\\
 &\leq26C_0\Delta^2\sum_{m=1}^{N}m
 \leq26C_0\delta^2.
\end{align*}
The second line uses \(m^{-1}\leq m\) and then
\(N(N+1)/2\leq N^2\) for \(N\geq1\).
The integral contribution from the same neighborhood is also of order
\(\delta^2\), since
\(\int_0^\delta r^2F(r)\,dr\leq C_0\delta^2/2\).

For \(R\geq1\), every point with \(\Delta|\nu|>R\) belongs
to a shell \(m>R/(\sqrt3\Delta)\). The second bound in
\eqref{v:eq:envolventes-termicas} gives
\begin{align*}
 \Delta^3\!\sum_{\Delta|\nu|>R}F(\Delta|\nu|)
 &\leq C_1e^{-dR/(2\sqrt3)}
 \Delta^3\sum_{m\geq1}(24m^2+2)e^{-d\Delta m/2}\\
 &\leq C_2e^{-dR/(2\sqrt3)},
\end{align*}
where \(C_2\) is independent of \(0<\Delta\leq1\). To verify
this uniformity, set \(q=e^{-d\Delta/2}\) and use
\[
 \sum_{m\geq1}q^m=\frac{q}{1-q},\qquad
 \sum_{m\geq1}m^2q^m=\frac{q(1+q)}{(1-q)^3}.
\]
The first identity is the geometric series, and the second follows by
applying \(q\,d/dq\) twice, which is valid for \(|q|<1\).
Moreover,
\(1-e^{-d\Delta/2}\geq\Delta(1-e^{-d/2})\) by concavity.
The factor \(\Delta^3\) therefore compensates for the denominators of
both expressions. The integral tail likewise tends to zero by the
exponential bound. For each \(\Delta>0\), the same estimates
prove convergence of the infinite sum.

On the compact annulus \(\delta\leq|k|\leq R\), the function
\(F(|k|)\) is continuous and bounded. Riemann sums of mesh
\(\Delta\) converge to its integral; the two boundary spheres
have measure zero. First let \(\Delta\to0\), then
\(\delta\to0\), and finally \(R\to\infty\), using the preceding
uniform bounds. The result is
\[
 \Delta^3\sum_{\nu\ne0}F(\Delta|\nu|)
 \longrightarrow \int_{\mathbb R^3}F(|k|)\,d^3k
 =4\pi\int_0^\infty r^2F(r)\,dr.
\]
Since \(2/L^3=2\Delta^3/(2\pi)^3\), the final radial coefficient
is \(2\cdot4\pi/(2\pi)^3=1/\pi^2\), proving
\eqref{v:eq:limite-modos-discretos}.
\end{proof}

Control at the origin is carried out over positive-energy excitations.
The constant mode remains fixed before the partition function is formed:
removing it afterward from an integral would not replace that condition on
the finite representation. The proof establishes the limit of the
specified realization and preserves the distinction between that realization and
the prior construction of its action and calendar parameters.


\subsection{Thermal moments and information}

\begin{lemma}[Bosonic moment]
\label{v:lem:momento-bosonico}
For every real \(s>1\),
\begin{equation}
 \int_0^\infty\frac{x^{s-1}}{e^x-1}\,dx
 =\Gamma(s)\zeta(s).
 \label{v:eq:momento-bosonico}
\end{equation}
\end{lemma}
\begin{proof}
For \(x>0\), \((e^x-1)^{-1}=\sum_{m\geq1}e^{-mx}\).
All summands are nonnegative; Tonelli's theorem allows the sum
and integral to be interchanged. The substitution \(y=mx\) gives
\(\int_0^\infty x^{s-1}e^{-mx}\,dx=m^{-s}\Gamma(s)\).
The sum converges for \(s>1\), yielding the stated equality.
\end{proof}

\begin{proposition}[Number, energy, and entropy densities]
\label{v:prop:densidades-termicas}
In the realization of Proposition~\ref{v:prop:ley-espectral},
\begin{align}
 n_\gamma&=\frac{2\zeta(3)}{\pi^2c^3(\beta\hbar)^3},\\
 u_\gamma&=\frac{\pi^2}{15c^3\hbar^3\beta^4},\\
 s_\gamma&=\frac{4\pi^2 k_B}{45c^3(\beta\hbar)^3}.
 \label{v:eq:densidades-termicas}
\end{align}
At the calendar temperature, these relations become
\begin{equation}
 n_\gamma\ell_P^3=\frac{2\zeta(3)}{\pi^2(\log3)^3},\quad
 \frac{u_\gamma\ell_P^3}{E_P}=\frac{\pi^2}{15(\log3)^4},\quad
 \frac{s_\gamma\ell_P^3}{k_B}=\frac{4\pi^2}{45(\log3)^3}.
 \label{v:eq:densidades-calendario}
\end{equation}
\end{proposition}

\begin{proof}
Integrating occupation against the mode density uses
\eqref{v:eq:momento-bosonico} with \(s=3\), and the energy integral
uses it with \(s=4\). This gives \(\Gamma(3)=2\) and
\(\Gamma(4)\zeta(4)=6\pi^4/90=\pi^4/15\).
For entropy, start from the partition product of the modes:
\[
 \frac{\log Z}{V}=-\frac{1}{\pi^2c^3}
 \int_0^\infty\omega^2
 \log(1-e^{-\beta\hbar\omega})\,d\omega
 =\frac{\pi^2}{45c^3(\beta\hbar)^3}.
\]
The last equality follows by integration by parts after the substitution
\(x=\beta\hbar\omega\). The boundary terms vanish:
\(x^3\log(1-e^{-x})\) tends to zero at both endpoints.
The Gibbs identity \(s_\gamma/k_B=\log Z/V+\beta u_\gamma\)
proves the third expression. Finally,
\(\beta_{108}\hbar=t_P\log3\) and \(\ell_P=ct_P\) give the three
calendar normalizations.
\end{proof}

\subsection{Photonic composition of Bendersky--Glaisher level \texorpdfstring{\(A_2\)}{A2}}
\label{v:subsec:bendersky-fotonica}

The number density in Proposition~\ref{v:prop:densidades-termicas}
has a second exact reading that retains its spectral provenance.
Let \(\mathcal Z_3\) denote the normalized scalar functional produced by
the tritic spectral reader and subsequently recognized as the zeta function
in the Archimedean realization. Let
\[
 H_k=\sum_{j=1}^{k}\frac1j,\qquad H_0=0,
\]
and let \(\mathsf B_j\) be the Bernoulli numbers under the convention
\(\mathsf B_1=-1/2\). To avoid any collision with lattice notation,
define the Bendersky--Glaisher levels by
\begin{equation}
 A_k^{\mathrm{BG}}
 =\exp\!\left(
   \frac{H_k\mathsf B_{k+1}}{k+1}-\mathcal Z_3'(-k)
 \right),\qquad k\in\mathbb N_0.
 \label{v:eq:bendersky-bg}
\end{equation}
The Bernoulli term fixes the normalization of the regularized product; it
does not introduce a physical constant. The coordinate
\(\pi_{\mathrm{HMT}}\) entering the functional equation has already
been expressed by the enriched APP--TRIT--TPK state before this reading.

\begin{proposition}[Photonic composition at level two]
\label{v:prop:bendersky-fotonica}
In the scalar normalization \(\mathcal Z_3=\zeta\),
\begin{equation}
 \log A_2^{\mathrm{BG}}
 =-\mathcal Z_3'(-2)
 =\frac{\mathcal Z_3(3)}{4\pi_{\mathrm{HMT}}^2}.
 \label{v:eq:bendersky-nivel-dos}
\end{equation}
Consequently, in the calendar thermal chart,
\begin{align}
 n_\gamma\ell_P^3
  &=\frac{8\log A_2^{\mathrm{BG}}}{(\log3)^3}
   =-\frac{8\mathcal Z_3'(-2)}{(\log3)^3},
   \label{v:eq:numero-fotones-bendersky}\\
 \frac{u_\gamma}{n_\gamma k_B\Theta_{108}}
  &=\frac{\pi_{\mathrm{HMT}}^2}{120\log A_2^{\mathrm{BG}}},
   \label{v:eq:energia-media-bendersky}\\
 \frac{s_\gamma}{n_\gamma k_B}
  &=\frac{\pi_{\mathrm{HMT}}^2}{90\log A_2^{\mathrm{BG}}}.
   \label{v:eq:entropia-media-bendersky}
\end{align}
\end{proposition}

\begin{proof}
The completed form satisfies
\[
 \pi_{\mathrm{HMT}}^{-s/2}\Gamma(s/2)\mathcal Z_3(s)
 =\pi_{\mathrm{HMT}}^{-(1-s)/2}
  \Gamma((1-s)/2)\mathcal Z_3(1-s).
\]
At \(s=-2\), the simple zero of \(\mathcal Z_3\) cancels the pole of
\(\Gamma(s/2)\). Comparing constant terms and using
\(\Gamma(3/2)=\sqrt{\pi_{\mathrm{HMT}}}/2\) after Archimedean
recognition gives
\[
 \mathcal Z_3'(-2)=-\frac{\mathcal Z_3(3)}{4\pi_{\mathrm{HMT}}^2}.
\]
Because \(\mathsf B_3=0\), Definition~\eqref{v:eq:bendersky-bg}
gives the first equality in~\eqref{v:eq:bendersky-nivel-dos}.
Substitution of
\(\mathcal Z_3(3)=4\pi_{\mathrm{HMT}}^2\log A_2^{\mathrm{BG}}\)
into~\eqref{v:eq:densidades-calendario} proves
\eqref{v:eq:numero-fotones-bendersky}. Finally,
\(k_B\Theta_{108}=E_P/\log3\); dividing the energy and entropy
densities by the number density proves
\eqref{v:eq:energia-media-bendersky} and
\eqref{v:eq:entropia-media-bendersky}.
\end{proof}

The evaluation \(A_2^{\mathrm{BG}}=1.0309167521973921\ldots\) is a
numerical test of the three identities, not generative input. The
Bendersky--Glaisher spectral development contains the complete hierarchy of
regularized derivatives and its correspondence with odd spectral values; the present section establishes the
reciprocal composition used specifically by dimensionless photon-number
density. It neither
redefines \(k_B\), identifies a density with a regularized operator, nor
alters any of the previously proved \(\zeta(3)\) densities.


\subsection{Thermal flux and the Stefan--Boltzmann constant}

\begin{proposition}[Isotropic flux and action normalization]
\label{v:prop:stefan}
The hemispherical flux of the isotropic realization is
\begin{equation}
 j_*(\Theta)=\sigma\Theta^4,\qquad
 \sigma=\frac{\pi^2 k_B^4}{60\hbar^3c^2}
 =\frac{2\pi^5k_B^4}{15h^3c^2}.
 \label{v:eq:stefan}
\end{equation}
At the calendar temperature,
\begin{equation}
 j_*(\Theta_{108})
 =\frac{2\pi^5}{15\,108^4(\log3)^4}
 \frac{h}{\ell_0^2t_0^2}.
 \label{v:eq:stefan-calendario}
\end{equation}
\end{proposition}
\begin{proof}
The isotropic energy per unit solid angle is
\(u_\gamma/(4\pi)\). Flux normal to a surface integrates
\(c\cos\theta\) over the hemisphere:
\[
 \frac{cu_\gamma}{4\pi}
 \int_0^{2\pi}\!d\phi\int_0^{\pi/2}\cos\theta\sin\theta\,d\theta
 =\frac{cu_\gamma}{4}.
\]
The formula for \(u_\gamma\), \(\beta^{-1}=k_B\Theta\), and
\(h=2\pi\hbar\) prove \eqref{v:eq:stefan}.
Substitution of \eqref{v:eq:energia-informacion-calendario} and
\(c=\ell_0/t_0\) gives \eqref{v:eq:stefan-calendario}.
\end{proof}

Within this realization, the constant \(\sigma\) is the coefficient
combining bosonic occupation, transverse-mode density, and
angular projection of the flux. The fourth-power dependence arises from the cubic
frequency moment together with the integration differential. The calendar
readout expresses that same flux through action, length,
time, and tritic information.

\newpage
\subsection{Spectral maxima and displacement law}

\begin{lemma}[Positive root for spectral maxima]
\label{v:lem:raiz-wien}
For every integer \(n>1\), the equation
\(n(1-e^{-x})=x\) has a unique solution \(x_n>0\).
The function \(f_n(x)=x^n/(e^x-1)\) has its unique positive maximum at
\(x_n\).
\end{lemma}
\begin{proof}
Let \(g_n(x)=n(1-e^{-x})-x\). We have
\(g_n(0)=0\), \(g_n'(x)=ne^{-x}-1\), and
\(g_n''(x)=-ne^{-x}<0\). The maximum of \(g_n\) is attained at
\(\log n\), where its value is \(n-1-\log n>0\), whereas
\(g_n(n)=-ne^{-n}<0\). Strict monotonicity on either side of the maximum
proves existence and uniqueness of the positive root. The sign of
\(f_n'(x)\) agrees with that of \(g_n(x)\); moreover, \(f_n\) tends
to zero at zero and infinity. This proves uniqueness of the maximum.
\end{proof}

\begin{proposition}[Frequency and wavelength readouts]
\label{v:prop:wien}
The maxima of radiance per frequency and per wavelength are
\begin{equation}
 \nu_{\max}=\frac{x_3 k_B\Theta}{h},\qquad
 \lambda_{\max}=\frac{hc}{x_5 k_B\Theta}.
 \label{v:eq:wien}
\end{equation}
In particular,
\begin{equation}
 \nu_{\max}(\Theta_{108})=
 \frac{x_3}{108t_0\log3},\qquad
 \lambda_{\max}(\Theta_{108})=
 \frac{108\log3}{x_5}\ell_0.
 \label{v:eq:wien-calendario}
\end{equation}
\end{proposition}
\begin{proof}
The conversion \(\omega=2\pi\nu\) and the isotropic angular factor give
\[
 B_\nu=\frac{2h\nu^3}{c^2}\frac1{e^{\beta h\nu}-1}.
\]
With \(\nu=c/\lambda\), the transformed density satisfies
\(B_\lambda=B_\nu|d\nu/d\lambda|\); hence
\[
 B_\lambda=\frac{2hc^2}{\lambda^5}
 \frac1{e^{\beta hc/\lambda}-1}.
\]
The respective changes of variable lead to \(f_3\) and \(f_5\).
The preceding lemma proves \eqref{v:eq:wien}; the calendar scale
gives \eqref{v:eq:wien-calendario}. The density Jacobian explains why
the two maxima correspond to distinct roots:
\(\nu_{\max}\lambda_{\max}=c x_3/x_5\).
\end{proof}

Wien displacement is structurally defined by the stationary point of
the chosen spectral density. A change of variable transforms
both the argument and the measure. This precision locates
each spectral constant in its generating operation and preserves its
physical meaning when passing between frequency and wavelength.
